\section{Related Work}
\label{frost-related-work}

\subsection{Prior Related Work}

We now review prior threshold schemes with a focus on Schnorr-based designs,
and split our review into robust and non-robust schemes. Robust schemes ensure that
so long as $t$ participants correctly follow the protocol, the protocol is
guaranteed to complete successfully, even if a subset of
participants (at most $n-t$) contribute malformed shares.
Conversely, designs that are not robust simply abort after detecting any
participant misbehaviour.

\textbf{Robust Threshold Schemes.} Stinson and Strobl~\cite{Stinson:2001:PSD}
present a threshold signature scheme
producing Schnorr signatures, using the modification of Pedersen's DKG presented
by Gennaro et al.~\cite{distributed-keygen} to generate
both the secret key $s$ during key generation as well as
the random nonce $k$ for each signing operation.
This construction
requires at minimum four rounds
for each signing operation (assuming no participant misbehaves): three rounds
to perform the DKG to obtain $k$,
and one
round to distribute signature shares and compute the group signature.
Each round requires participants to send values to every other participant.

Gennaro et al.~\cite{distributed-keygen-revisited} present a threshold
Schnorr signature protocol that uses a modification of Pedersen's
DKG~\cite{pedersen:DKG} to generate both $s$ during key generation
and the random nonce $k$ for signing operations. However, their construction
requires \emph{all} $n$ signers to participate in signing, while the adversary
is allowed to control up to the given threshold number of participants.
Recall from Section~\ref{background:distributed-keygen} that Pedersen's DKG
requires two rounds; this construction requires an additional round for signing
operations when all participants are equally trusted. Each round requires that
all participants send values to all other participants.
The authors also discuss an optimization that leverages
a \emph{signature aggregator} role, an entity trusted to
gather signatures from each participant, perform validation, and publish the
resulting signature, a role we also adopt in our work.
In their optimized variant, participants can perform
Pedersen's DKG to generate multiple $k$ values in a pre-processing stage independently
of performing signing operations.
In this variant, to compute $\ell$ signatures, signers first perform two
rounds of $\ell$ parallel executions of Pedersen's DKG,
thereby generating $\ell$
random nonces. The signers can then store these pre-processed values to later perform
$\ell$ single-round signing operations.

Our work builds upon the key generation stage of Gennaro et
al.~\cite{distributed-keygen-revisited}; we use a variant of Pedersen's DKG for
key generation with a requirement
that in the case of misbehaviour, the protocol aborts and the cause investigated
out of band.
However, FROST \emph{does not} perform a DKG
during signing operations as is done in both of the above schemes, but instead
make use of additive secret sharing and share conversion.
Consequently, FROST trades off robustness for more efficient signing
operations,
such that a misbehaving participant can cause the signing operation to abort.
However, such a tradeoff is practical to many real-world settings.

Further, because FROST does not provide robustness, FROST is secure so long as
the adversary controls fewer than the threshold $t$ participants,
an improvement over robust designs, which can at best
provide security for $t \le n/2$~\cite{distributed-keygen}.

\textbf{Non-Robust Threshold Schemes.} FROST is not unique in trading off
favouring increased network efficiency over robustness.
Gennaro and Goldfeder~\cite{gennaro-goldfeder-ecdsa} present a threshold ECDSA
scheme that similarly requires aborting the protocol in the case of
participant misbehaviour. Their signing construction uses a two-round DKG to generate
the nonce required for the ECDSA signature, leveraging additive-to-multiplicative share
conversion. This DKG has been also applied in a Schnorr threshold
scheme context to generate the random nonce for more efficient distributed key
generation operations~\cite{kzen-schnorr} in combination with threshold Schnorr
signing operations~\cite{Stinson:2001:PSD}.
In later work~\cite{one-round-ecdsa}, Gennaro and Goldfeder define an optimization
to a single-round ECDSA signing operation with a
preprocessing stage, which assumes the protocol will abort in the case of
failure or participant misbehaviour.
Their end-to-end protocol with identifiable aborts has eight network rounds,
six of which require broadcasting to all other signing participants,
and two of which require performing pairwise multiplicative-to-additive
share conversion protocols. Further, while the protocol can be optimized into a
preprocessing phase, the choice of the signing coalition must be determined at
the time of preprocessing.
FROST defines a more efficient preprocessing
phase as secret nonces can be generated in a distributed manner in the
preprocessing phase entirely
non-interactively. Further,
participants can ``mix'' preprocessed values across different signing
coalitions, as FROST requires that the choice for the signing coalition be made
only during the signing stage.

Recent work by Damg\aa{}rd et al.~\cite{fast-threshold-ecdsa} define an efficient
threshold ECDSA construction that similarly requires aborting in the case of
misbehaviour. Their design relies on generating a blinding factor
$d + m \cdot e$ such that
where $d$ and $e$ are $2t$ secret sharings of zero, such
that the entire binding factor evaluates to zero when all signing parties
are honest and agree on $m$. This approach is similar to FROST in that signature
shares are bound to the message and to the set of signing parties.
However, the security of their
scheme requires the majority of participants to be honest, and $n \geq 2t+1$.
Further, their scheme requires all $n$ participants take part in signing
operations, where the threshold $t$ is simply a security parameter.

Similarly to FROST,
Abidin, Aly, and Mustafa~\cite{schnorr-thresh-additive} present a design for
authentication between devices,
and use additive secret sharing to generate the nonce for Schnorr
signatures in a threshold setting, a technique
also used by FROST.
However, the authors do not consider the
Drijvers attack and consequently their design is similarly limited to
restricted levels of parallelism. Further, their design does
not include validity checks for responses submitted by participants when
generating signatures and consequently does not detect nor identify
misbehaving participants.

FROST improves upon prior work in Schnorr threshold schemes by
providing a single-round signing variant with a preprocessing stage that is
agnostic to the choice of the signing coalition. Further, the number of signing
participants in FROST is required to be simply some $t\leq n$,
while remaining secure against the
Drijvers attack and misbehaving participants who do not correctly follow the
protocol.









\subsection{Subsequent Work}

Since the publication of FROST, many Schnorr-related schemes have been proposed.
We now review this subsequent work.



%%%%%%%%%%%%%%%%%%%%%%%%
Sparkle related work



\paragraph{Threshold Schnorr Signatures.}
Closest to the design of $\sparkleds$ is the MSDL scheme presented
by Boneh, Drijvers, and Neven~\cite{BonehDN18},
the three-round MuSig scheme by Maxwell, Poelstra, Seurin, and
Wuille~\cite{MaxwellPSW19},
and the $2$Schnorr scheme by Nicolosi, Krohn, Dodis, and Mazi{\`{e}}res~\cite{NicolosiKDM03}.
However, MSDL and MuSig consider only
the multi-signature setting ($n$-out-of-$n$),
and $2$Schnorr considers only the $2$-out-of-$2$ setting.
Note that when proving the security of multi-signatures, there is only one honest signer.

Stinson and Strobl~\cite{StinsonR01} propose a threshold Schnorr signature scheme secure in the
random oracle model under the discrete logarithm assumption. However,
their scheme requires performing a three-round distributed key generation protocol (DKG)~\cite{GennaroJKR07} to generate the nonce for
each signature, which adds considerable network overhead: at a minimum, it requires participants to perform four rounds in total.
Furthermore, the proof of security assumes only a static adversary.

Komlo and Goldberg \cite{KomloG20} present a two-round threshold Schnorr signature FROST.
Unlike prior threshold Schnorr schemes in the literature~\cite{GennaroJKR03}, FROST is
secure against a concurrent adversary and is not susceptible to ROS attacks~\cite{BenhamoudaLLOR22}.
%that is allowed to open simultaneous signing sessions at once.
%In other words, FROST is secure even against ROS attacks~\cite{BenhamoudaLLOR21}.
%Additionally, FROST can be seen as a non-interactive threshold signature because the first round of signing can be preprocessed before the signing set and message are known.
FROST2, proposed by Crites, Komlo, and Maller~\cite{CritesKM21} and refined by Bellare, Tessaro, and Zhu~\cite{BellareTZ22} and BCKMTZ~\cite{BellareCKMTZ22}, is an optimized version of FROST that reduces the number of exponentiations required for signing from $\thresh$ to one.
A further optimization of FROST2, called FROST3, is introduced by RRJSS~\cite{RuffingRJSS22} and proven secure with a DKG by Chu, Gerhart, Ruffing, and Schr{\"{o}}der~\cite{ChuGRS23}.
(See Table~\ref{table:efficiency} for a comparison of efficiency.)
However,
all three of FROST, FROST2, FROST3 are only proven statically secure, under the algebraic one-more discrete logarithm assumption (AOMDL) \cite{BellareCKMTZ22}.
While $\sparkleds$ adds an additional round of communication,
it requires only a single exponentiation per signer and is proven statically secure under standard assumptions, which are also criteria of interest in~\cite{NIST22, NIST23}.

%(2) it only requires a constant number of exponentiations per signer, and (3) it is provably secure against a concurrent adversary in the adaptive setting.
%CK- I took out (2) and (3) since this statement makes it sounds like neither FROST/FROST2 can achieve either 2 or 3. I think the most important and salient point here is sparkle is secure under standard assumptions.
% Liz - The first statement is too weak so I removed it

Lindell presents a three-round threshold Schnorr signature~\cite{Lindell22} with the goal of
defining a scheme that is both secure against ROS attacks and
secure in the random oracle model under the discrete logarithm assumption only.
Security is modeled in the the universally composable framework (UC)~\cite{Canetti01} and therefore captures a concurrent adversary.
However, the security proof assumes the adversary is static, and no claims are made regarding adaptive security.
%Compared to the scheme by Lindell,
$\sparkleds$ similarly relies on only the ROM and DL assumption for static security, but does not require the use of
online-extractable zero-knowledge proofs~\cite{Fischlin05}, and hence is both
significantly more efficient and a simpler design.
\begin{comment}
  % CK - not sure if we actually need this
In estimating the cost of the online extractor, which in practice requires the
Fischlin transform~\cite{Fischlin05}, we estimate that to achieve 128-bit
security, generating a Fischlin proof requires the number of commitments $r$ to equal 16 and the length of
the zero vector $b$ to equal $16$.
\end{comment}

Concurrent to this work,
Makriyannis~\cite{Makriyannis22} defines a commit-reveal threshold
Schnorr signature scheme similar to $\sparkleds$, called Classic S., and proves security with
respect to an idealized notion of threshold signatures by CGGMP21~\cite{CanettiGGMP21}.
They consider adaptive security but employ the guessing argument that incurs exponential tightness loss relative to the number of parties.

\begin{figure}[t]
	\centering
%		\resizebox{\textwidth}{!}{
		\begin{tabular}{ c | c | c | c | c | c | c | c | c }
                   & \multicolumn{5}{c|}{} & \multicolumn{2}{c|}{} \\
      Scheme        & \multicolumn{5}{c|}{$\Sign$} &
      \multicolumn{2}{c|}{$\Combine$} \\
             \midrule
			%\cmidrule{2-3} \cmidrule{5-8} \cmidrule{9-10} \cmidrule{11-12}
      &  \multicolumn{3}{c|}{{\small~Performance~}} &
      \multicolumn{2}{c|}{{\small~Bandwidth~}} &
      \multicolumn{2}{c|}{{\small~Performance~}} \\
      &  rounds & exp & $\Hash$ &  $~\Gr~$ & $~\F$  & exp & $\Hash$
		\\ \midrule
      FROST \cite{KomloG20}        		  	&    $2$ 	& $t+2~$	    & $t+1$       &  $2$	  & $1$	    &  $t$	& $t$		 \\
      FROST2 \cite{BellareCKMTZ22}      		&    $2$ 	& $3$	        & $2$         &  $2$	  & $1$	    &  $1$	& $1$	 \\
      FROST3 \cite{RuffingRJSS22}      		&    $2$ 	& $3$	        & $2$         &  $2$	  & $1$	    &  $1$	& $1$	 \\
      Classic S.~\cite{Makriyannis22}	        					       &    $3$ 	& $1$	        & $1$         &  $1$	  & $2$	    &  $0$	& $0$		 \\
      %Zero S.~\cite{Makriyannis22}		& $3$	& 			&		&		&			& \\
      Lindell22 \cite{Lindell22}~        &    $3$ 	& $11t + 1$   & $18t+10^7$ 	&  $11$		& $25$		&  $0$  & $0$		\\
      \hline
      \textbf{$\sparkleds$}         					       &    $3$ 	& $1$	        & $t + 2$         &  $1$	  & $2$	    &  $0$	& $0$		 \\ \midrule
% SpeedyBSig
%      \multicolumn{11}{l}{\textbf{Blind signatures}}   \\ \midrule
%      %Abe~\cite{Abe01} & $X$ & $X$ & $X$   & - &  $X$ & $X$    &  $X$    & $X$  &  $X$      & $X$ \\
%      TZ22~\cite{TessaroZ22} & $1$ & $2$ & $0$   & - &  $2$ & $11$    &  $4$    & $7$  &  $-$      & $4$ \\
%      Our $\speedybsig$               & $1$ & $2$ & $0$   & - &  $2$ & $10$  &  $3$     & $6$  &  $-$   & $3$ \\
      \bottomrule
	\end{tabular}
%}
	\caption{Efficiency of Two- and Three-Round Threshold Schnorr Signature Schemes.
	All output a standard Schnorr signature.
	We only compare schemes that are secure against ROS attacks~\cite{DrijversEFKLNS19,BenhamoudaLLOR22}.
		The number of network rounds between participants is given in the rounds column.
  exp stands for the number of group exponentiations.
  The total number of group and field elements sent by each signer is denoted by $\Gr$ and $\F$, respectively.
  $\Hash$ denotes the total number of hashing operations performed.
  The cost of signature verification is identical for each scheme, and is simply the cost of verifying a single Schnorr signature.
  Estimates for Lindell22 are made with respect to a 128-bit security level for
  Fischlin~\cite{Fischlin05}, where $r=8$ is the number of commitments for a
  Fischlin proof and the length of the zero vector is $b=16$, such that
  $b\cdot r=128$.
	}\label{table:efficiency}
	\hrulefill
	\vspace{-0.5cm}
\end{figure}


\paragraph{Adaptive Security of Threshold Signatures.}
While adaptive security for threshold schemes is a well-known topic in the
literature,
no proof of adaptive security without exponential tightness loss exists for
a threshold Schnorr signature scheme that is secure against ROS attacks.
%We now review several alternative signature schemes, for comparison.

Generalized techniques for transforming statically secure threshold schemes into adaptively secure schemes have been defined in the literature~\cite{CanettiGJKR99, JareckiL00, LysyanskayaP01}.
However, these techniques either require the reduction to guess the corrupted parties ahead of time,
require secure erasures, or
introduce prohibitive performance overhead by using a robust distributed key generation mechanism (DKG) for nonce generation or heavyweight tools, such as Paillier encryption.

Almansa, Damg{\aa}rd, and Nielsen~\cite{AlmansaDN06} present a threshold RSA
scheme with proactive and adaptive security, but these results do not translate to the discrete logarithm setting.
Libert, Joye, and Yung~\cite{LibertJY16} propose a variant of the threshold
BLS~\cite{Boldyreva03} scheme that is adaptively secure.
However, it is incompatible with single-party BLS verification, an
often-critical goal for threshold schemes in practice.
Bacho and Loss~\cite{BachoL22} demonstrate the adaptive security of threshold
BLS in Type I bilinear groups directly in the AGM from the $\thresh\hbox{-}\omdl$ assumption.
Their reduction is tight.
Interestingly, they demonstrate that the $\thresh\hbox{-}\omdl$ assumption is the minimum assumption under which threshold BLS can be proven adaptively secure.
\liz{Check this, because I think we say something different in the impossibility paper.}

A recent result by Das and Ren~\cite{DasR23} shows adaptive security of threshold BLS in Type III bilinear groups under the co-computational Diffie-Hellman assumption (co-CDH) and decisional Diffie-Hellman assumption (DDH) in the ROM.

Makriyannis~\cite{Makriyannis22} introduces a threshold Schnorr signature scheme, Zero S., which relies on Pedersen commitments, Fischlin proofs of knowledge~\cite{Fischlin05},
and an interactive protocol to identify misbehaving parties.
Zero S. is proven adaptively secure in a straight-line manner with a reduction to DL but requires secure erasure of secret state.

Twinkle, by BLTWZ~\cite{BachoLTWZ23}, is a family of three-round threshold signature schemes.
The first, efficient construction is a distributed Chaum-Pedersen proof of equality of discrete logarithms $\pk = g^{\sk}$ and $\Hash(m)^{\sk}$~\cite{ChaumP92} (vs. a proof of knowledge of the discrete logarithm of $\pk = g^{\sk}$, as in Schnorr signatures).
Thus, signatures are Chaum-Pedersen outputs $(\Hash(m)^{\sk}, c, z)$, where $(c, z)$ is a Schnorr signature.
This scheme is proven fully adaptively secure under a one-more variant of the computational Diffie-Hellman assumption (OMCDH) in the ROM.
The second, matrix construction is proven fully adaptively secure under the DDH assumption in the ROM.
The techniques used for simulation do not allow pre-processing: the message and signing set must be given in the first round.
Furthermore, neither scheme is compatible with Schnorr signatures.
To date, ours is the only proof of adaptive security for any threshold Schnorr signature scheme.
